\section{Open questions}

\begin{enumerate}
\item \textbf{End-to-end composition (paper's C5).} Does the $\alpha$-QPE / RFPE
acceleration survive when composed with a real molecular VQE loop, and what is
the empirical shot / wall-time advantage vs $\alpha{=}0$ at fixed
chemical-accuracy target? \textit{Basis:} the paper reports the VQE and
$\alpha$-QPE ingredients separately (C1 and C2/C3) but never runs the
composition; constants (particle count, prior widening between VQE outer steps,
rejection-filter acceptance at deep $M$) could erase the analytical advantage.
\textit{Next steps:} drive PennyLane VQE on H$_2$, LiH, BeH$_2$ with an
$\alpha$-QPE expectation subroutine at $\alpha\in\{0,0.25,0.5,0.75,1\}$;
measure shots and coherent time to chemical accuracy; report cost-Pareto vs
$\alpha$; check if the analytical scaling law holds under composition.

\item \textbf{$\alpha$-QPE + classical shadows.} Can $\alpha$-QPE be combined
with post-2020 classical-shadow tomography (Huang--Kueng--Preskill,
derandomised / locally-biased variants) to further cut per-expectation shot
cost, especially on many-Pauli-term Hamiltonians of larger molecules?
\textit{Basis:} the two accelerations attack orthogonal axes (per-Pauli-term
sampling vs coherent depth per phase estimate) and are candidate-multiplicative.
\textit{Next steps:} hybrid subroutine --- shadows for the variance-dominant
Pauli terms, $\alpha$-QPE ($\alpha \in [0.5, 1]$) for the residual
high-variance terms; benchmark H$_2$O/STO-3G and N$_2$/6-31G shot counts
vs plain shadows and plain $\alpha$-QPE at matched target variance.

\item \textbf{Noise robustness of the acceleration.} How does the $\alpha$-QPE
acceleration degrade under realistic hardware noise (T1/T2, gate infidelity,
readout error), and is there a noise-dependent optimal $\alpha^\star < 1$?
\textit{Basis:} the paper's Fig.~5 and this replication are noiseless. Larger
$\alpha$ buys precision by increasing coherent depth $M_k = 1/\sigma_k^\alpha$,
which is the regime decoherence should first kill.
\textit{Next steps:} rerun RFPE with a depolarising / amplitude-damping channel
of strength $\lambda$ per unit coherent time; sweep $(\alpha, \lambda)$; find
$\alpha^\star(\lambda)$ maximising Bayes-risk shrinkage rate; compare against
super-conducting (T1 $\sim100\,\mu$s) and trapped-ion (T1 $\sim$s) regimes.

\item \textbf{Integration with adaptive ansatz methods.} Does $\alpha$-VQE
compose cleanly with adaptive-ansatz methods (ADAPT-VQE, ROTOSELECT,
qubit-ADAPT) that add operators one at a time by gradient screening?
\textit{Basis:} adaptive ans\"atze want cheap low-precision gradients for
selection and expensive high-precision evaluations for refinement --- the
variable-precision regime where $\alpha$ as a knob should be strongest, but
the paper studies only fixed-precision $\alpha$-QPE in a fixed ansatz.
\textit{Next steps:} implement ADAPT-VQE with an $\alpha$-scheduled
subroutine ($\alpha_{\rm screen}\in[0, 0.25]$ during ranking,
$\alpha_{\rm refine}\in[0.75, 1]$ for the retained-operator VQE step);
benchmark on H$_2$, LiH, H$_4$ chain vs uniform-$\alpha$ $\alpha$-VQE and
plain ADAPT-VQE.

\item \textbf{Sensitivity to classical-optimizer choice.} How sensitive is the
observed $\alpha$-QPE acceleration to outer-loop optimiser choice
(Adam / SPSA / QNG / L-BFGS), and does the optimal $\alpha$ shift with
optimiser noise-tolerance? \textit{Basis:} this replication used Adam
(first-order, comparatively noise-robust) with a noiseless subroutine;
optimisers that consume gradient variance differently (SPSA noise-native;
L-BFGS not) should prefer different $\alpha$.
\textit{Next steps:} sweep $\{$Adam, SPSA, QNG, L-BFGS$\}\times\alpha$ on H$_2$
and LiH; report shots-to-chemical-accuracy heatmap; check if the ordering
$C_{\alpha{=}1} < C_{\alpha\in(0,1)} < C_{\alpha{=}0}$ is preserved or if some
optimisers invert the advantage.
\end{enumerate}
